\chapter{Prezentacja warstwy użytkowej projektu}
\section{Zrzuty ekranu aplikacji}

W tej części wstawiono zrzuty ekranu. Pokazują one najważniejsze widoki aplikacji i odpowiadają strukturze programu \texttt{HotelManagementApp}.

\screenplaceholder{01_ekran_glowny.eps}{Ekran główny aplikacji z kartami statystyk i opisem modułów}

\screenplaceholder{02_goscie.eps}{Widok zarządzania gośćmi hotelowymi}

\screenplaceholder{03_pokoje.eps}{Widok pokoi hotelowych z typem pokoju i statusem}

\screenplaceholder{04_rezerwacje.eps}{Widok tworzenia rezerwacji dla wybranego gościa i pokoju}

\screenplaceholder{05_platnosci.eps}{Widok płatności przypisanych do rezerwacji}

\screenplaceholder{06_pracownicy.eps}{Widok pracowników hotelu}

\screenplaceholder{07_uslugi.eps}{Widok usług z możliwością dodawania i edycji}


\section{Opis działania najważniejszych funkcjonalności}

Najważniejszą częścią aplikacji jest możliwość zarządzania rezerwacjami hotelowymi. Użytkownik wybiera gościa, pokój i termin pobytu. Program sprawdza, czy pokój jest dostępny, a następnie oblicza cenę na podstawie typu pokoju i liczby nocy.

Drugą ważną funkcją jest obsługa płatności. Użytkownik wybiera rezerwację, wpisuje kwotę i zapisuje płatność. Dane są widoczne w tabeli płatności oraz wpływają na statystyki na panelu głównym.

Kolejna funkcja to moduł usług. Użytkownik może dodać nową usługę, na przykład parking albo śniadanie, oraz edytować istniejącą usługę. Formularz został przygotowany tak, aby jedna sekcja służyła do dodawania i edycji.

\section{Prezentacja interfejsu użytkownika}

Interfejs został podzielony na trzy główne części. Na górze znajduje się nagłówek z nazwą systemu i przyciskiem odświeżenia danych. W środku znajduje się obszar roboczy z zakładkami. Komunikaty o błędach i wykonanych operacjach są pokazywane w formularzach lub w aktywnej sekcji programu, zależnie od wykonywanej operacji.

Panel główny pokazuje najważniejsze statystyki: aktywne rezerwacje, liczbę gości, dostępne pokoje, liczbę usług i wartość płatności. Dzięki temu użytkownik po uruchomieniu programu od razu widzi aktualny stan systemu.

\section{Przykładowe scenariusze działania aplikacji}

Pierwszy scenariusz dotyczy dodania gościa. Użytkownik przechodzi do widoku gości, wpisuje imię, nazwisko, telefon, adres e-mail i numer dokumentu, a następnie naciska przycisk dodania. Jeżeli dane są poprawne, gość pojawia się w tabeli.

Drugi scenariusz dotyczy utworzenia rezerwacji. Użytkownik wybiera gościa, pokój, pracownika oraz datę zameldowania i wymeldowania. Po zapisaniu rezerwacja pojawia się na liście. Jeżeli pokój jest zajęty w podanym terminie, aplikacja pokazuje komunikat o braku dostępności.

Trzeci scenariusz dotyczy edycji usługi. Użytkownik przechodzi do widoku usług, wybiera usługę z tabeli, zmienia nazwę lub cenę i naciska przycisk zapisu. System zapisuje zmiany w bazie danych.
